\settrack{tracktwo}
% VOICE: 1st-person personal (Bruce) — "I" not "Bruce". Personal/investigative voice.

\chapter{The Departure}
\label{record:departure}

\begin{quote}\small\textit{%
``In some sort of crude sense which no vulgarity, no humor, no overstatement can quite extinguish, the physicists have known sin; and this is a knowledge which they cannot lose.''%
}\par\raggedleft --- Oppenheimer, MIT lecture (1947)\end{quote}

\vspace{0.5cm}

% DMS MVP import from staging/raw/ — not final prose

\srcnote{HEALER-RECONSTRUCTION.md | staging/raw/pos07-the-departure.md | 2026-02-15 | We Already Did Mate --- Joy article as disclosure, close textual comparison}

\section*{The Revelation}
\label{record:dep-the-revelation}

\srcnote{V9 spatial grounding | staging/interview-content-prepared.md | 2026-03-15 | 15th and Grant scene, 8-year deferred understanding}

Fall 2004, Corvallis. I was sitting in the living room of my house on 15th and Grant, mid-morning. My girlfriend was away. The kids were at school. The house was quiet in the way houses are when you're about to ask a question you've been assembling for weeks.

By then I had pieced together enough of the implications to frighten myself. Dual-use technology. Nuclear launch codes\ldots{} A potential arms race that would make the Cold War look like a border dispute. I had been running the scenarios in my head, and every one of them ended badly.

Healer told me that he and his colleagues wanted to ``step into the light'' and say ``we did this.'' They feared that by the time their work was declassified, they would be too old, not listened to, or dead. This book is a move towards that outcome.

I asked: ``Why not just tell the truth now, in a way that no one would notice?''

David looked at me with that expression I'd come to know --- the one that meant I was standing on the answer and couldn't see it.

``We already did, mate.''

I didn't understand. I filed it away the way you file things you know are important but can't yet parse. It sat there for eight years --- a locked box with no key. I carried it through the departure, through the silence, through years of solitary research where I chased every thread David had given me without knowing what picture they formed.

% SPIRAL-REPEAT: "Why the Future Doesn't Need Us" — key reference, also in pos06-the-secret.tex, pos27-extension.tex, timeline.tex, abstracts.tex
He was talking about Bill Joy's article ``Why the Future Doesn't Need Us,'' published in \textit{Wired} magazine in April 2000.\footcite{joy2000} I did not figure this out until 2012, when my friend and researcher Mark R.\ from Eugene found it. Mark had been helping me research for about four months when he emailed me a link. I opened it at my desk in Corvallis, facing the fireplace.

I read Joy's essay straight through. Electric shock. Tingling. \textit{This is the thing I missed.} Joy names the same scientists. Describes the same fears. Uses the same word --- relinquishment. Published April 1, 2000. April Fools' Day. The biggest truth disguised as the biggest joke.

I sat motionless for a long time. Then I went for a fifty-mile bike ride --- two and a half hours in the saddle --- because my body needed to move while my mind caught up. When I got back I began the intensive research that eventually became this book.

Eight years. ``We already did, mate'' --- and the answer had been sitting in \textit{Wired}'s archive the entire time.

\section*{The Reread}
\label{record:dep-the-reread}

Under one reading --- the one most people give it --- ``Why the Future Doesn't Need Us'' is a brilliant technologist's warning about the future. Under another reading, it is something else entirely. The second reading occurred to me in 2012. It has not left me since.

\vspace{0.5cm}

What follows is either the most revealing close reading in this book or the most damning exhibit of pattern-matching. I cannot tell you which.

\begin{quote}
\textit{A ten-point close textual comparison between Joy's 2000 \textit{Wired} essay and my reconstruction --- phrase matches, named circles, Kauffman's citation, Hillis's calm, the April Fools' publication date --- is collected in \autoref{app:joy-ten-point}.}
\end{quote}

\vspace{1cm}\noindent\rule{0.5\textwidth}{0.5pt}\vspace{0.5cm}

\section*{The Departure}

\srcnote{Autobiography of an Accidental Conspirator by Bruce.txt | staging/raw/pos03-the-mentor.md | 2026-02-15 | Departure narrative --- August-September 2006}

In September 2006, I was told my services would not be required. My children could not be protected --- the divorce, the custody arrangements, the impossibility of shielding them from what was coming. David delivered this without cruelty. It was an operational assessment, not a rejection of me as a person. The full significance of that assessment --- what they needed protection from --- belongs to a chapter not yet written.

The project that proposed to protect all of humanity could not protect one man's daughters.

David and I parted. His decision, not mine. He said a fond goodbye but didn't drag it out. He was not a man who lingered.

I was no less dedicated. But I didn't know what I had been dedicated to. That realization would take years to arrive --- and when it did, it came gradually.

Under Possibility~A, the ``vetting rejection'' was a convenient exit from a relationship built on deception --- David leaving before the story fell apart. Under Possibility~B, David may have been involved in sensitive work, but framing children as security vulnerabilities imports the language of espionage onto what may have been a simpler situation. Under Possibility~C, a man who had been selected for a role in something extraordinary was told he couldn't continue, and spent the next four years locked behind glass --- able to think about what had happened but unable to speak it.

\chapterreturn
